\subsection{Every ramified local section of the full cubic pro-three group}

\begin{statement}{E.61}{Proved}{The nine full local cusp families and the abelian inertia image of metabelian sections}
Let $K=\mathbb Q(\zeta_3)$, let $C$ be the smooth cubic
$X^3+Y^3+Z^3=0$, and remove its nine points on the coordinate
axes to obtain $U$. Take $(1:-1:0)$ as origin on $C$.
Use the marked $[3]$-cover identification of E.50 and E.58, so that
\[
 \Pi=\pi_1(U_{\overline K})^{(3)}=N
 \quad\hbox{is free pro-three of rank ten}.
 \tag{E61.1}
\]
Fix a finite place $v\nmid3$, put $q=|k_v|$ and
$h=v_3(q-1)\ge1$, and use an integral nonzero tangent at the
base cusp. The resulting genuine geometric action of $G_{K_v}$
on $\Pi$ is unramified. Fix compatible roots defining a tame
three-primary inertia generator $t$ and arithmetic Frobenius $F$.
For each of the nine cusps, choose an integral tangent and an
inertia-invariant geometric path from the base tangent. Write
$I_i=\langle d_i\rangle\simeq\mathbb Z_3(1)$ for its oriented
geometric inertia and $F_i=z_iF$ for its actual tangential
Frobenius lift. In particular, $F_id_iF_i^{-1}=d_i^q$.

The continuous sections with nontrivial geometric inertia,
modulo conjugation by $\Pi$, form exactly the following space:
\[
 \coprod_{i=1}^{9}
       (\mathbb Z_3\setminus\{0\})\times
             \mathbb Z_3/(q-1)\mathbb Z_3.
 \tag{E61.2}
\]
The actual representative attached to $(i,a,\bar b)$ is
\[
       t\longmapsto d_i^a t,\qquad
       F\longmapsto d_i^bF_i.
 \tag{E61.3}
\]
Changing $b$ by $(1-q)k$ is precisely conjugation by $d_i^k$.
The cusp and the nonzero exponent $a$ are unchanged by any
geometric conjugation between representatives. Formula (E61.2)
is a homeomorphism for the quotient topology on continuous
cocycles and the indicated three-adic and finite topologies.

Put $M=\Pi/\Pi''$, with its induced genuine arithmetic action,
and $V=\Pi^{\rm ab}$. The image of the continuous map taking
a local $M$-section class to its abelian inertia is exactly
\[
 \operatorname{im}\bigl(
   \operatorname{Sect}(M\rtimes G_{K_v})/M
       \longrightarrow V\bigr)
       =\bigcup_{i=1}^{9}\mathbb Z_3[d_i].
 \tag{E61.14}
\]
The map uses the fixed tame generator $t$ above; geometric
conjugation acts trivially on its value in $V$. Each vector
on the right is realized by a full cuspidal section, including
the zero vector. This statement computes the abelian inertia
image and retains the further commutator and Frobenius fibers.

At $K_v=\mathbb Q_2(\zeta_3)$, let
$\Pi\twoheadrightarrow R$ be the actual arithmetic-equivariant quotient
of E.58. Number the cusps by $(i,j)\in\{0,1,2\}^2$ so that
their images are $c^{e_{ij}}$. In the normalized local
coordinates of E.60, the image of every full ramified section
has
\[
   (A,B)=(i,j),\qquad
   (h_{kl})=a e_{ij},\qquad u=v=0,
       \quad a\ne0.                                  \tag{E61.4}
\]
The continuous $R$-section
\[
       \nu=c^{\,e_{00}+e_{01}+V},\qquad z=\eta^2
 \tag{E61.5}
\]
has no lift to a continuous section with geometric kernel
$M$, and hence has no full $\Pi$-lift, even after full
geometric conjugation. The map $\Pi\twoheadrightarrow R$
factors through $M\twoheadrightarrow R$, since the commutator
subgroup of $R$ lies in its displayed abelian module.
Its inertia has an actual word preimage
$c(ycy^{-1})[y^3,c]$.
The obstruction is already detected by its image
$e_{00}+e_{01}$ in
\[
 \bigoplus_{i,j}\mathbb F_3 e_{ij}\big/
             \mathbb F_3\sum_{i,j}e_{ij}
          \ \lhook\joinrel\longrightarrow\ \Pi^{\rm ab}/3.
 \tag{E61.6}
\]
In the associated finite-cover proof the detecting cover has
degree $19683$, has $59049$ cusps, and its compactification has
genus $59050$. These integers and all twenty-seven possible
cuspidal-line comparisons in (E61.6) are checked by the retained
program. This settles the entire ramified local source at these
good places and gives an explicit reverse-lifting failure for
the computed quotient. The global and unramified source
questions remain at their stated maps in E.23.
The manuscript anchor is \S51.16, folios 733v--736v,
scans 172--178, formulas (371)--(379), Part E(c),(d).
\end{statement}

\begin{proof}
\emph{The actual specialization and local source.}
The Fermat equation is smooth away from three: its three
partial derivatives have no common projective zero there.
The nine cusp sections are disjoint since the third roots of
unity remain distinct. With the displayed origin they are
$C[3]$; the identity over $K$ extends over the good integral
model because $[3]$ is etale there and all nine sections
remain distinct. Thus the special fibre has genus one, nine
rational punctures, and rational three-torsion.

Apply \cite[XIII, Corollary 2.9 and construction 2.10]{E_SGA1}
to this smooth proper model with its disjoint cusp divisor,
using finite three-groups. The specialization functor identifies
every finite three-group torsor on the two geometric fibres,
compatibly with its fibre functor and the cusp inertia
\cite[XIII, Lemma 2.11]{E_SGA1}. It induces a continuous
isomorphism of the two geometric pro-three groups. The
integral tangential fibre functor makes this isomorphism
equivariant for $G_{K_v}\twoheadrightarrow G_{k_v}$ and makes
the based action unramified. We therefore work with the actual
special-fibre arithmetic extension
\[
           \Pi\rtimes G_{k_v},\qquad G_{k_v}=\widehat{\mathbb Z}.
 \tag{E61.7}
\]
The genus-one, nine-puncture presentation follows as well
from \cite[XIII, Corollary 2.12]{E_SGA1}; eliminating any one
cusp generator gives ten free pro-three generators.

On $\Pi^{\rm ab}/3$, Frobenius preserves the exact sequence
whose subspace is the eight-dimensional cusp space and whose
quotient is $C[3]$. It acts as the identity on each, so its
matrix is unipotent. The kernel of
$\operatorname{Aut}(\Pi)\to\operatorname{Aut}(\Pi^{\rm ab}/3)$
is pro-three, as seen from the characteristic exponent-three
central series. Hence the closed Frobenius action group is
pro-three. Every cocycle together with this action is a
homomorphism into a pro-three semidirect product and therefore
factors through $G_{K_v}^{(3)}$.

The latter quotient is the tame group
\[
       D=\langle t,F\mid FtF^{-1}=t^q\rangle_{(3)}.
 \tag{E61.8}
\]
Indeed wild inertia is pro-residue-characteristic, its image
in a pro-three group is trivial, and the prime-to-three parts
of tame inertia and of the unramified quotient have trivial
image. The usual tame relation survives with this exact $q$;
it defines the semidirect product
$\mathbb Z_3\rtimes\mathbb Z_3$ because $q\in1+3\mathbb Z_3$.
The same argument supplies the universal map in the reverse
direction from the tame quotient to every such pro-three
image. Thus a pair $(\nu,z)\in\Pi^2$ defines a continuous
section exactly when
\[
                  zF(\nu)z^{-1}=\nu^q.                \tag{E61.9}
\]

\emph{Recognition of the nonzero inertia in the full group.}
We give the pro-three, finite-field form of the argument in
\cite[Theorem 30]{E_StixCuspidal}, including the cover and its
comparison map. The cited paper works with full profinite
groups in characteristic zero; the following argument verifies
the precise variant used here.

Every cusp group $I_i$ is a free pro-three factor of $\Pi$.
Any two cusp generators occur simultaneously in a free basis,
by eliminating a third cusp in the surface presentation.
Theorem B' of \cite{A_HerfortRibes1985} applies to the extension-
closed class of finite three-groups: a nontrivial element of
a free factor has its centralizer in that factor, and the
factor has trivial intersection with any different conjugate
of itself. Different cusp factors have disjoint nonzero
conjugates as well: an equality
$d_i^a=g d_j^b g^{-1}$ in $\Pi$ gives
$a[d_i]=b[d_j]$ in the free integral abelianization, whose two
cusp vectors are independent. This forces $a=b=0$ for $i\ne j$.
These facts also show that every procyclic subgroup meeting
a cusp group nontrivially is contained in that group: its
generator centralizes the nonzero intersection. Every nonzero
procyclic subgroup is isomorphic to $\mathbb Z_3$, since the
free pro-three group is torsion-free (Theorem A' of the same
source applied to its free pro-three cyclic factors).

Take a solution of (E61.9) with $\nu\ne1$, and put
$I=\overline{\langle\nu\rangle}\simeq\mathbb Z_3$ and
$\psi=zF$. Then $\psi$ normalizes $I$ by its $q$-power action.
Let $W$ be the union of all conjugates of all nine cusp groups.
It is compact, being the image of
$\Pi\times\coprod_i I_i$ under conjugation.
Suppose $I$ is contained in no cusp group. The preceding
centralizer calculation gives $I\cap W=\{1\}$.
The compact set $I\setminus I^3$ is disjoint from $W$.
Compact separation in a profinite group supplies an open
normal subgroup for which their images are disjoint.
Intersect its finitely many automorphic images to obtain a
characteristic open subgroup $N_0$ with the same property;
there are finitely many subgroups of any fixed index because
$\Pi$ is finitely generated. In particular,
$I\cap N_0\subset I^3$.

The subgroup $H=N_0I$ is open and $\psi$-stable. Its exact
quotient map
\[
 H\longrightarrow I/(I\cap N_0)
           \longrightarrow I/I^3\simeq\mathbb F_3
 \tag{E61.10}
\]
is nonzero on $I$ and kills every cusp subgroup of $H$.
For otherwise some element of $W\cap H$ would have, modulo
$N_0$, a representative in $I\setminus I^3$, contrary to the
choice of $N_0$.

The continuous map $\widehat{\mathbb Z}\to\Pi\rtimes
\widehat{\mathbb Z}$ taking $1$ to $\psi$ is a section:
its composition with the projection is the identity.
Consequently $H\rtimes\langle\psi\rangle$ is an open
arithmetic subgroup and defines an actual geometrically
connected finite etale cover $U'\to U_{k_v}$ over $k_v$.
Its geometric group is $H$. More explicitly, any finite
three-cover of its geometric fibre has a three-group Galois
closure over $U_{\overline{k_v}}$, by the finite wreath-product
construction; thus taking pro-three completion loses none
of the covering subgroups of $H$.
Let $C'$ be the smooth compactification of $U'$.
Filling the cusp inertia gives the continuous surjection
\[
 H\twoheadrightarrow\pi_1(C'_{\overline{k_v}})^{(3)}
       \twoheadrightarrow T_3\operatorname{Jac}(C').
 \tag{E61.11}
\]
The last arrow is abelianization, with its canonical
identification from \cite[III, Proposition 9.1]{E_MilneAV2008}.
It is independent of the Abel--Jacobi base point after
abelianization and is Frobenius-equivariant. The kernel of
the first arrow is precisely the closed normal subgroup
generated by the puncture inertia, since an etale cover of
$U'$ extends etale to $C'$ exactly when its inertia is trivial.
Hence (E61.10) factors through (E61.11). The image of $I$ in
the free $\mathbb Z_3$-module $T_3\operatorname{Jac}(C')$
is nonzero modulo three.

Conjugation by $\psi$ acts as arithmetic Frobenius on this
Tate module and by multiplication by $q$ on that nonzero
image. Thus $q$ would be an eigenvalue. The actual Frobenius
eigenvalues on the Tate module of an abelian variety over
$\mathbb F_q$ have complex absolute value $\sqrt q$
\cite[II, Theorem 1.1(b), Lemmas 1.2--1.3]{E_MilneAV2008}.
Since $q>1$, the asserted eigenvalue is impossible. This
proves that $I$ lies in a unique conjugate of a cusp group.
Notice that the weight argument has been applied to the
finite-field Jacobian in (E61.11), with its constructed
Frobenius action.

\emph{Evaluation of the full ramified section space.}
Conjugate the unique containing cusp group to the fixed
$I_i$. Then $\nu=d_i^a$ for one nonzero $a\in\mathbb Z_3$.
The element $\psi$ preserves $I_i$, since it preserves its
nonzero subgroup $I$. Both $\psi$ and the chosen $F_i$
normalize $I_i$ and project to $F$. Their ratio lies in the
geometric normalizer, which is $I_i$. It follows that
$\psi=d_i^bF_i$ for a unique $b\in\mathbb Z_3$.
Conversely these elements satisfy (E61.9) directly, and
(E61.8) proves continuity and existence of the entire section.

A geometric conjugation between two such representatives
must take the unique containing cusp group to the other one.
The disjointness already proved forces the same cusp and
places the conjugator in $I_i$. Conjugation by $d_i^k$ fixes
$a$ and sends $b$ to $b+(1-q)k$. This proves (E61.2)--(E61.3),
including every equivalence relation. Continuity follows from
three-adic exponentiation. For the inverse, the primitive
integral cusp vectors in $\Pi^{\rm ab}$ determine the cusp
and the nonzero exponent continuously. On each clopen set
$v_3(a)=r$, the parameter space is compact and its continuous
bijection to the corresponding Hausdorff conjugacy quotient
is a homeomorphism. These sets exhaust the nonzero-inertia
space; their clopen property is detected by the same
abelianization. This proves the topological assertion.

\emph{The exact abelian inertia image of metabelian sections.}
Let a continuous $M$-section have tame coordinates
$(\bar\nu,\bar z)$, and let $w$ be the image of $\bar\nu$
in $V$. Its abelian tame relation is $F(w)=qw$. We prove
that $w$ lies on one of the nine integral cusp lines.
Suppose it does not. The compact primitive annulus
$\mathbb Z_3w\setminus3\mathbb Z_3w$ is disjoint from
$\bigcup_i\mathbb Z_3[d_i]$. Indeed each $[d_i]$ is primitive
in the free lattice $V$; a nonzero intersection of its rational
line with $\mathbb Q_3w$ would put $w$ on its integral line.
Choose $n$ so that these compact sets remain disjoint in
$V/3^nV$. This also gives
$\mathbb Z_3w\cap3^nV\subset3\mathbb Z_3w$.

Let $H$ be the inverse image of $3^nV+\mathbb Z_3w$ under
$\Pi\twoheadrightarrow V$. It is open, contains $\Pi'$, and
is preserved by $zF$ for any lift $z\in\Pi$ of $\bar z$.
The exact quotient character
\[
 H\longrightarrow
   (3^nV+\mathbb Z_3w)/3^nV
       \longrightarrow\mathbb Z_3w/3\mathbb Z_3w
 \tag{E61.15}
\]
detects the image of $\bar\nu$ and kills every cusp group of
$H$, by the chosen separation. Because $\Pi'\subset H$,
we have the actual closed-subgroup inclusion
$\Pi''\subset[H,H]$. There is therefore a continuous quotient
map
\[
 H/\Pi''\longrightarrow H^{\rm ab}
       \longrightarrow T_3\operatorname{Jac}(C'),
 \tag{E61.16}
\]
where $C'$ is the compactification of the finite-field cover
defined by $H\rtimes\langle zF\rangle$ exactly as in
(E61.11). The relation of the proposed section holds modulo
$\Pi''$, so becomes the exact eigenvalue-$q$ relation in
(E61.16). Character (E61.15) factors through this same Tate
module and detects a nonzero image modulo three. The
finite-field Weil bound again excludes that eigenvalue.
This proves the inclusion in (E61.14). Every vector on its
right occurs in (E61.3), including $a=0$ for the zero vector,
which proves equality with the stated image.

\emph{The actual quotient comparison and its failed lift.}
The images of the marked cusp words
$x^iy^j c y^{-j}x^{-i}$ in $R$ are $c^{e_{ij}}$.
The full tangential Frobenius at this cusp has handle
coordinates $(i,j)$ in the normalization of E.60: in the
ambient two-handle abelianization, Frobenius is $-2I$, so
$pF(p)^{-1}$ for $p=x^iy^j$ has coordinates $(3i,3j)$.
Inner conjugation in $\Pi=N$ changes these normalized
coordinates by multiples of three. The nonzero cusp inertia
becomes exactly $c^{a e_{ij}}$, proving (E61.4).
The additional finite Frobenius coordinates depend on the
chosen actual paths and tangents. This statement does not
assign them from a reference automorphism of the ambient
free group: the remaining inner parameter described in E.58
can change those origins while acting trivially on $R$.

For (E61.5) the moment data are
\[
       A=0,\quad B=2,\quad h_0=2,\quad I=0,\quad J_0=1,
              \quad u=0,\quad v=1.
 \tag{E61.12}
\]
The exact E.60 defect is
$4[(Ah_0-I-3u)U+(Bh_0-J_0-3v)V]=0$.
The retained program also evaluates the actual eleven-module
Frobenius and the group law, obtaining
$z\varphi(\nu)z^{-1}\nu^{-4}=1$.
It evaluates the displayed word preimage using the same group
law. Thus this is a continuous section of the actual quotient.
Its image under the isomorphism $R^{\rm ab}\simeq\Pi^{\rm ab}$
of E.58 has cusp vector $e_{00}+e_{01}$. For every cusp label
$k$ and every $a\in\mathbb F_3$, the vector
$e_{00}+e_{01}-a e_k$ has unequal coordinates and so cannot
belong to the diagonal line. This checks all twenty-seven
possibilities in (E61.6). Geometric conjugation acts trivially
on this abelianization. The exact metabelian image (E61.14)
therefore forbids every metabelian lift of (E61.5), and the
full classification gives the same obstruction for $\Pi$.

There is a finite geometric cover detecting exactly this
obstruction. Put $\Phi\Pi=\Pi^3[\Pi,\Pi]$ and let $w$ be
the nonzero vector (E61.6). The inverse image $H$ of the line
$\langle w\rangle$ in $\Pi/\Phi\Pi\simeq\mathbb F_3^{10}$
has index $3^9=19683$. Each cusp vector is outside this line,
so each cusp has ramification index three in the cover;
the number of lifted cusps is $9\cdot3^8=59049$.
The punctured Euler characteristic is multiplied by $3^9$,
and hence the compact genus is
\[
       1+\frac{9\cdot3^9-59049}{2}=59050.
 \tag{E61.13}
\]
For a proposed metabelian lift, choose any representative
$z\in\Pi$ of its Frobenius coordinate. Its $zF$ stabilizes this exact
$H$, by (E61.9). The quotient character
$H\to H/\Phi\Pi=\langle w\rangle$ kills all its cusps and
detects its inertia. Here $\Pi'\subset\Phi\Pi\subset H$
gives $\Pi''\subset[H,H]$, so its relation descends along
the literal map (E61.16). Thus (E61.11) gives the same impossible
Frobenius eigenvalue four on the compact Jacobian of this
specific cover. The finite certificate verifies the line
separation, the cover indices and Euler characteristic, and
the cusp parameter formulas modulo $3,9,27,81$.
\end{proof}
